% ---------------------------------------------------------------------------
%  VIII. Conclusion + acknowledgment
%  Rule:    answer the research question of section 01 in its own terms,
%           and introduce nothing that has not already been established.
%           The acknowledgment must stay free of institution names while the
%           anonymisation policy is in force.
% ---------------------------------------------------------------------------

\section{Conclusion}

We asked whether the command surface of a closed Android--\ros{} service robot can be recovered
without the manufacturer's help, and what evidence that recovery demands. The seven-phase
pipeline produced a working surface on a commercial reception platform and an edge deployment
that replaces the vendor cloud for teleoperation, navigation, speech and visitor interaction. The
vendor application supplied protocol names and undocumented preconditions (H1); the exposed
broker could not carry commands at all (H2). Most importantly, the gap between transport and
execution (H4) governed the exercise: the addressing scheme we expected to matter (H3) explained
none of our outcomes, while the precondition state explained all of them. Establishing that
required correcting our own measurement twice, our clearest evidence for H4: the instrument is
as easily fooled as the integrator. Next steps: a second manufacturer, a measured task comparison
in place of our capability table, and a language-model planner grounded in these primitives
rather than left to rediscover the interface itself.


% DOUBLE-BLIND: acknowledgments are omitted from the submitted version, since
% naming a laboratory or a funder identifies the authors. Restore for the
% camera-ready version:
%
%   \section*{Acknowledgment}
%   We thank our robotics laboratory for robot access and for support during
%   field testing.
